\documentclass{amsart}
\usepackage{amsmath,amssymb,amsthm,hyperref}
\usepackage{algorithm}
\usepackage{algpseudocode}

\newtheorem{theorem}{Theorem}
\newtheorem{lemma}{Lemma}
\newtheorem{proposition}{Proposition}
\newtheorem{corollary}{Corollary}
\theoremstyle{definition}
\newtheorem{definition}{Definition}
\newtheorem{remark}{Remark}
\newtheorem{assumption}{Assumption}

\newcommand{\Chat}{\widehat{C}}
\newcommand{\Lhat}{\widehat{L}}
\newcommand{\Nhat}{\widehat{N}}
\newcommand{\E}{\mathbb{E}}
\renewcommand{\Pr}{\mathbb{P}}

\begin{document}

\title{Measurement, estimation, and provable mapping quality in heterogeneous computing}
\maketitle

\section{Overview}\label{sec:overview}

The problem asks us to (a) define computable runtime descriptors $L_i(\tau)$ and
$N_{ij}(\tau)$, (b) produce completion-time estimators $\Chat$ with controlled
accuracy, and (c) prove that minimizing the \emph{estimated} objective yields a
mapping provably close to the optimum of the \emph{exact} objective.

The mathematical heart of the matter is the following implication, which we
prove in full generality:

\begin{quote}
\emph{If the estimator is accurate in a multiplicative (relative-error) sense
uniformly over all candidate mapping/schedule pairs, then any
mapping/schedule that exactly minimizes the estimated objective is, under the
exact objective, within a factor $\tfrac{1+\varepsilon}{1-\varepsilon}$ of
optimal, for a broad and explicitly characterized class of objectives (makespan
under precedence and weighted average completion time among them).}
\end{quote}

We then close the loop with concrete measurement and estimation procedures.
The crucial subtlety, which an earlier version of this manuscript elided, is
that the descriptors and hence the completion times \emph{depend on the chosen
mapping and schedule} through contention, queueing, and precedence. We
therefore (i) make the completion time an explicit function of the schedule via
a \emph{structural demand/load model}, (ii) estimate the schedule-independent
\emph{demand} parameters of that model from runtime probes, and (iii) propagate
the estimation error through the model to obtain uniform multiplicative accuracy
\emph{simultaneously over all schedules}, with an explicit sample-complexity
bound that carries an honest accounting of estimator bias, sample dependence,
and temporal drift.

\medskip
\noindent\textbf{Roadmap.} Section~\ref{sec:model} fixes the schedule model and
the structural completion-time model. Section~\ref{sec:objective} defines
scale-coherent objectives and verifies the canonical ones, now for the
schedule-induced finish times. Section~\ref{sec:measure} defines and computes
the descriptors. Section~\ref{sec:estimate} estimates the structural demand
parameters and propagates the error to a \emph{schedule-uniform} multiplicative
accuracy guarantee, accounting for bias, dependence, and drift.
Section~\ref{sec:transfer} proves the transfer theorem and its robust and
end-to-end corollaries.

\section{Model: schedules and structural completion times}\label{sec:model}

\subsection{Schedulable units, mappings, and schedules}

\begin{definition}[Units and precedence]\label{def:units}
Let $U=T\cup S$ be the finite set of \emph{schedulable units}: tasks
$T=\{t_1,\dots,t_k\}$ and the (finite) set $S$ of all subtasks obtained by
decomposition. The data/precedence relations form a directed acyclic graph
(DAG) $G=(U,E)$: an edge $(u,v)\in E$ means $v$ may not start before $u$
finishes and (if $u,v$ are placed on different machines) before the data
produced by $u$ has been transferred to $v$'s machine. Each unit $u$ carries an
input/output \emph{data size} $d_{uv}\ge0$ for each edge $(u,v)\in E$.
\end{definition}

\begin{definition}[Mapping and schedule]\label{def:sched}
Let $M=\{m_1,\dots,m_n\}$. A \emph{mapping} is a function $x\colon U\to M$;
write $x_u=i$ if $u$ is assigned to $m_i$. A \emph{schedule} $\sigma$ for a
mapping $x$ assigns to each unit $u$ a start time $\sigma(u)\ge0$ on machine
$x_u$ and to each cross-machine edge a transfer interval, such that
\begin{enumerate}
\item[(S1)] no two units mapped to the same machine overlap in time;
\item[(S2)] every precedence/data constraint of $G$ is respected: $v$ starts no
earlier than $u$'s finish plus the relevant transfer time for each $(u,v)\in E$.
\end{enumerate}
A \emph{configuration} is a pair $\pi=(x,\sigma)$ satisfying (S1)--(S2) together
with the problem's hard placement constraints (co-location, exclusivity,
availability). Let $\Pi$ be the (finite, nonempty) set of admissible
configurations. We assume $\Pi\neq\emptyset$ (e.g.\ any greedy list schedule on
a feasible mapping is admissible).
\end{definition}

The finiteness of $\Pi$ as a \emph{decision} set deserves a word: start times
range over a continuum, but for any fixed mapping the relevant schedules are
determined up to combinatorial choice (the relative order of units on each
machine and the routing of transfers), and the optimal start times are then the
earliest feasible ones given that order. Thus the optimization ranges over a
finite set of \emph{orderings}, each inducing a unique left-shifted schedule;
we identify $\Pi$ with this finite set. This is the standard reduction used for
shop and DAG scheduling.

\subsection{Structural completion-time model}

The completion time of a unit is \emph{not} a primitive constant: it depends on
the machine, on the unit's intrinsic resource \emph{demand}, and on the
\emph{load} the configuration places on that machine and its links. We make
this explicit.

\begin{definition}[Demand parameters]\label{def:demand}
Each unit $u$ has a \emph{demand vector} $d_u\in\mathbb{R}^{q}_{\ge0}$
(e.g.\ floating-point operations, memory footprint, I/O volume) intrinsic to the
unit and independent of placement. Each cross-machine edge $(u,v)$ carries a
transferred data size $d_{uv}\ge0$. Collect all of these into the
\emph{demand profile} $\mathbf d=(\,(d_u)_{u},(d_{uv})_{(u,v)\in E}\,)$, a fixed
but a priori unknown vector in $\mathbb{R}^Q_{\ge0}$.
\end{definition}

\begin{definition}[Structural completion-time / finish-time model]\label{def:struct}
There is a known family of \emph{per-unit completion functions}
\[
C_{ui}\;=\;g_i\big(d_u;\,L_i,\,N_{\cdot i}\big)\;>\;0,
\]
the time to run $u$ on $m_i$ \emph{in isolation given a load/status profile},
where $g_i$ is monotone and \emph{positively homogeneous of degree one in
$d_u$ when the profile is held fixed} (doubling the work doubles the time),
which is the standard linear-throughput model
$g_i(d_u;L_i,N)=\langle \kappa_i(L_i),\,d_u\rangle$ with rate vector
$\kappa_i(L_i)>0$ depending on the machine's status. Given a configuration
$\pi=(x,\sigma)$, the descriptors $L_i,N$ are themselves induced by $\pi$
(queues, utilizations, effective bandwidths are functions of who is placed
where and when); write $L_i(\pi),N(\pi)$. The \emph{realized completion time}
of unit $u$ in configuration $\pi$ is
\[
C_u(\pi)\;:=\;g_{x_u}\big(d_u;\,L_{x_u}(\pi),\,N(\pi)\big),
\]
and the \emph{finish time} $f_u(\pi)$ is the schedule-induced time at which $u$
completes, i.e.\ $f_u(\pi)=\sigma(u)+C_u(\pi)$ with $\sigma$ the earliest
feasible left-shifted schedule of the configuration (so (S1)--(S2) and transfer
delays $\lambda_{ij}+d_{uv}/\beta_{ij}$ are respected).
\end{definition}

The key structural fact we will use is that $C_u(\pi)$ — and through it the
finish times and every objective below — is \emph{positively homogeneous of
degree one in the demand profile $\mathbf d$}, with the configuration held
fixed.

\begin{lemma}[Homogeneity of the finish times in demand]\label{lem:homog}
Fix a configuration $\pi\in\Pi$. For every $c>0$, replacing $\mathbf d$ by
$c\,\mathbf d$ replaces every realized completion time $C_u(\pi)$ by
$c\,C_u(\pi)$, every transfer delay's data-dependent part by $c$ times itself,
and hence every left-shifted finish time $f_u(\pi)$ by $c\,f_u(\pi)$, provided
the latency part of transfers and the schedule \emph{ordering} are held fixed.
\end{lemma}

\begin{proof}
By Definition~\ref{def:struct}, $C_u(\pi)=\langle\kappa_{x_u}(L_{x_u}(\pi)),d_u\rangle$
with the profile $L,N$ \emph{frozen at $\pi$}; this is linear in $d_u$, so
scaling $d_u\mapsto c d_u$ scales $C_u(\pi)\mapsto cC_u(\pi)$. A cross-edge
transfer time is $\lambda_{ij}+d_{uv}/\beta_{ij}$; its data-dependent part
$d_{uv}/\beta_{ij}$ scales by $c$. The earliest left-shifted finish time of a
fixed ordering is a finite maximum of sums of completion and transfer times
along chains of the DAG (the standard recursion
$f_u=\max(\text{machine predecessor finish},\ \max_{(v,u)\in E}(f_v+\text{transfer}))+C_u$).
If \emph{all} additive constituents that depend on $\mathbf d$ scale by $c$ and
the latency constants are treated as part of the fixed schedule overhead,
each $f_u$ is a max/sum of $c$-scaled terms, hence scales by $c$. (When pure
latencies $\lambda_{ij}$ are present and nonzero they are an additive
$\mathbf d$-independent constant; the homogeneity then holds for the
work-dependent component, and Remark~\ref{rem:latency} records the resulting
additive correction in the transfer theorem.)
\end{proof}

\section{Objectives and scale-coherence}\label{sec:objective}

We now phrase the objective directly on configurations through the realized
finish times, and isolate the structural property that drives the transfer
theorem.

\begin{definition}[Objective on configurations]\label{def:obj}
An \emph{objective} is a function $F\colon\mathbb{R}^Q_{\ge0}\times\Pi\to
\mathbb{R}_{\ge0}$, $F(\mathbf d,\pi)$, the cost of running configuration $\pi$
under demand profile $\mathbf d$. The canonical objectives are
\begin{align}
\text{(makespan)}\quad &F_{\max}(\mathbf d,\pi):=\max_{u\in U} f_u(\pi),
\label{eq:makespan}\\
\text{(weighted completion)}\quad &F_{w}(\mathbf d,\pi):=\sum_{u\in U} w_u\, f_u(\pi),
\label{eq:wsum}
\end{align}
with $w_u\ge0$, where $f_u(\pi)$ are the realized finish times of
Definition~\ref{def:struct}. Makespan is the time the last unit finishes —
correctly accounting, through $f_u$, for overlap, waiting, and precedence
(unlike the naive sum of per-unit times).
\end{definition}

\begin{definition}[Demand-scale-coherence]\label{def:coherent}
An objective $F$ is \emph{demand-scale-coherent} if it is built from the finish
times $f_u(\pi)$ by a map that is (i) \emph{monotone} (increasing each $f_u$
weakly increases $F$) and (ii) \emph{positively homogeneous of degree one in the
finish-time vector}: $F$ evaluated at finish times $(c f_u)_u$ equals
$c\cdot F$ evaluated at $(f_u)_u$, for all $c>0$. Combined with
Lemma~\ref{lem:homog}, this yields the operative property: for fixed $\pi$ and
all $c>0$,
\begin{equation}\label{eq:scale}
F(c\,\mathbf d,\pi)=c\,F(\mathbf d,\pi),
\end{equation}
and monotonicity in $\mathbf d$: $\mathbf d\le\mathbf d'$ (entrywise) implies
$F(\mathbf d,\pi)\le F(\mathbf d',\pi)$ for every $\pi$.
\end{definition}

\begin{lemma}[The canonical objectives are demand-scale-coherent]\label{lem:canonical}
$F_{\max}$ and $F_w$ are demand-scale-coherent.
\end{lemma}

\begin{proof}
By Lemma~\ref{lem:homog}, scaling $\mathbf d\mapsto c\mathbf d$ with $\pi$ fixed
scales every $f_u(\pi)\mapsto c\,f_u(\pi)$. Then
$F_{\max}=\max_u f_u$ scales as $\max_u c f_u=c\max_u f_u$, and
$F_w=\sum_u w_u f_u$ scales as $\sum_u w_u c f_u=c\sum_u w_u f_u$; both are
positively homogeneous of degree one in $(f_u)_u$, giving \eqref{eq:scale}.
Monotonicity: increasing any demand entry weakly increases every realized
completion time (each $\kappa_i>0$) and hence, through the left-shift
recursion, weakly increases every finish time; $\max$ and nonnegatively
weighted sums are monotone in their arguments. Hence both objectives are
monotone in $\mathbf d$.
\end{proof}

\begin{remark}\label{rem:classes}
Demand-scale-coherence holds for any monotone, degree-one-homogeneous aggregate
of finish times: $\ell_p$ norms of the finish-time vector ($p\ge1$), maximum
weighted finish time $\max_u w_u f_u$, makespan over per-machine loads, and
convex combinations thereof. The transfer theorem below uses \emph{only}
Definition~\ref{def:coherent} and so applies to all of these verbatim.
\end{remark}

\section{Measurement: computing descriptors and demands (part (a))}\label{sec:measure}

\begin{definition}[Descriptors]\label{def:descriptors}
At sampling instant $\tau$ each machine $m_i$ exports
$L_i(\tau)=(q_i(\tau),\rho^{\mathrm{cpu}}_i(\tau),\rho^{\mathrm{mem}}_i(\tau),
a_i(\tau),\theta_i)\in\mathbb{R}^5$ (queue length, CPU/memory utilization,
availability, nominal throughput), and each link exports
$N_{ij}(\tau)=(\beta_{ij}(\tau),\lambda_{ij}(\tau))\in\mathbb{R}^2_{\ge0}$
(effective bandwidth, latency). These determine the rate vectors $\kappa_i(L_i)$
and the transfer parameters used in Definition~\ref{def:struct}.
\end{definition}

\begin{algorithm}[h]
\caption{ComputeDescriptors}\label{alg:descr}
\begin{algorithmic}[1]
\Require window length $r\in\mathbb N$; raw streams $L_i(\tau_1),\dots,L_i(\tau_r)$, $N_{ij}(\tau_1),\dots,N_{ij}(\tau_r)$.
\Ensure averaged descriptors $\bar L_i,\bar N_{ij}$.
\For{$i=1$ \textbf{to} $n$}
  \State $\bar L_i\gets\frac1r\sum_{p=1}^r L_i(\tau_p)$
  \For{$j\ne i$} \State $\bar N_{ij}\gets\frac1r\sum_{p=1}^r N_{ij}(\tau_p)$ \EndFor
\EndFor
\State \Return $(\bar L_i)_i,(\bar N_{ij})_{ij}$
\end{algorithmic}
\end{algorithm}

\noindent\textbf{Correctness, termination, complexity.} A fixed number of
additions/divisions: it computes exactly the empirical means, terminates, and
runs in $O(rn^2)$ time, $O(n^2)$ space.

The descriptors feed the estimator in two ways: they instantiate the rate
vectors $\kappa_i(\bar L_i)$ and transfer rates $(\bar\beta_{ij},\bar\lambda_{ij})$
of the structural model, and they are used to estimate the demand profile
$\mathbf d$, as we describe next.

\section{Estimation: demands $\Rightarrow$ schedule-uniform accuracy (part (b))}\label{sec:estimate}

The earlier version probed the completion time $C_{ui}$ \emph{directly}, which is
unsound: $C_{ui}$ depends on the configuration, so a single probe corresponds to
a single load condition and cannot certify accuracy under the (different) load
of an optimal schedule, and one cannot probe a configuration without deploying
it. We instead estimate the \emph{schedule-independent} demand parameters
$\mathbf d$ and \emph{compute} $\Chat$ for every configuration via the known
structural model. This converts a per-configuration accuracy problem (with
$|\Pi|$ exponentially large) into a per-demand-entry accuracy problem (with $Q$
entries), and propagates a uniform relative error to \emph{all} configurations
at once.

\begin{assumption}[Probing model]\label{ass:probe}
For each demand coordinate $\ell\in\{1,\dots,Q\}$ (i.e.\ each $d_u$ component or
each $d_{uv}$), we obtain measured probes $Y^{(1)}_\ell,\dots,Y^{(r_\ell)}_\ell$
of $d_\ell$ — e.g.\ instruction/FLOP counters, allocated-byte counters, and
transferred-byte counters, read under a lightly loaded calibration run or
directly from hardware performance counters that measure \emph{work}, not
elapsed time. We assume:
\begin{enumerate}
\item[(M1$'$)] \emph{Bounded bias.} $|\E[\bar Y_\ell]-d_\ell|\le\beta\,d_\ell$
for a known $\beta\in[0,1)$, where $\bar Y_\ell=\frac1{r_\ell}\sum_p Y^{(p)}_\ell$.
(Counters of \emph{work performed} are nearly unbiased; $\beta$ budgets any
residual systematic error. Setting $\beta=0$ recovers exact unbiasedness.)
\item[(M2)] \emph{Boundedness.} $0\le Y^{(p)}_\ell\le b_\ell$ a.s., $b_\ell>0$ known.
\item[(M3)] \emph{Positivity.} $d_\ell\ge a_\ell>0$ known.
\item[(M4)] \emph{Effective independence.} The probes form a stationary sequence
whose autocorrelations $\rho_k$ are summable; with
$\tau_{\mathrm{ac}}:=1+2\sum_{k\ge1}|\rho_k|<\infty$ the effective sample size is
$r_\ell^{\mathrm{eff}}:=r_\ell/\tau_{\mathrm{ac}}$. (For i.i.d.\ micro-probes,
$\tau_{\mathrm{ac}}=1$.)
\end{enumerate}
\end{assumption}

\begin{proposition}[Per-coordinate multiplicative accuracy]\label{prop:samples}
Fix $\varepsilon\in(0,1)$ with $\varepsilon>\beta$, and $\delta\in(0,1)$. Put
$\varepsilon_0:=\varepsilon-\beta>0$. Under
Assumption~\ref{ass:probe}, if
\begin{equation}\label{eq:rbound}
r_\ell^{\mathrm{eff}}\ \ge\ \frac{b_\ell^2}{2\,a_\ell^2\,\varepsilon_0^2}\,
\ln\frac2\delta
\qquad\Big(\text{i.e. } r_\ell\ge \tau_{\mathrm{ac}}\,\frac{b_\ell^2}{2a_\ell^2\varepsilon_0^2}\ln\tfrac2\delta\Big),
\end{equation}
then $\Pr\big[(1-\varepsilon)d_\ell\le \bar Y_\ell\le(1+\varepsilon)d_\ell\big]\ge1-\delta$.
\end{proposition}

\begin{proof}
Write $\mu_\ell:=\E[\bar Y_\ell]$. By (M2) and the effective-sample-size
correction in (M4), a Hoeffding bound for stationary bounded sequences (or
Azuma--Hoeffding applied to the martingale of partial deviations, or simply
Hoeffding applied to $r_\ell^{\mathrm{eff}}$ independent blocks) gives, for any
$s>0$,
\[
\Pr\big[|\bar Y_\ell-\mu_\ell|\ge s\big]\le 2\exp\!\Big(-\frac{2 r_\ell^{\mathrm{eff}} s^2}{b_\ell^2}\Big).
\]
Take $s=\varepsilon_0 a_\ell$. The right side is $\le\delta$ exactly when
\eqref{eq:rbound} holds. On the complementary event,
$|\bar Y_\ell-\mu_\ell|<\varepsilon_0 a_\ell\le\varepsilon_0 d_\ell$ by (M3).
Combining with the bias bound (M1$'$), $|\mu_\ell-d_\ell|\le\beta d_\ell$, the
triangle inequality gives
\[
|\bar Y_\ell-d_\ell|\ \le\ |\bar Y_\ell-\mu_\ell|+|\mu_\ell-d_\ell|
\ <\ \varepsilon_0 d_\ell+\beta d_\ell\ =\ \varepsilon\, d_\ell,
\]
i.e.\ $(1-\varepsilon)d_\ell<\bar Y_\ell<(1+\varepsilon)d_\ell$.
\end{proof}

\begin{corollary}[Uniform demand accuracy]\label{cor:uniform}
Define $\widehat{\mathbf d}:=(\bar Y_\ell)_{\ell=1}^Q$. If for each coordinate
$r_\ell^{\mathrm{eff}}\ge\frac{b_\ell^2}{2a_\ell^2\varepsilon_0^2}\ln\frac{2Q}{\delta}$,
then with probability $\ge1-\delta$,
\begin{equation}\label{eq:Deps}
(1-\varepsilon)\,d_\ell\ \le\ \widehat d_\ell\ \le\ (1+\varepsilon)\,d_\ell
\qquad\text{for all }\ell=1,\dots,Q.
\end{equation}
\end{corollary}

\begin{proof}
Apply Proposition~\ref{prop:samples} with $\delta\mapsto\delta/Q$ to each
coordinate and take a union bound over the $Q$ coordinates.
\end{proof}

We call event \eqref{eq:Deps} the $\varepsilon$-accuracy event $E_\varepsilon$.
The expected relative error per coordinate is
$\E|\bar Y_\ell-d_\ell|/d_\ell\le\beta+\sqrt{\operatorname{Var}(\bar Y_\ell)}/d_\ell
\le\beta+b_\ell/(2a_\ell\sqrt{r_\ell^{\mathrm{eff}}})=\beta+O(1/\sqrt r)$,
so for $\beta\to0$ it is $O(1/\sqrt r)$.

\begin{lemma}[Error propagation: demands $\Rightarrow$ completion times]\label{lem:propagate}
Suppose $E_\varepsilon$ of \eqref{eq:Deps} holds. Build the estimated
completion times by the \emph{same} structural model evaluated at the estimated
demands and the measured descriptors:
$\Chat_u(\pi):=g_{x_u}(\widehat d_u;\bar L_{x_u},\bar N)$ and estimated transfer
work $\widehat d_{uv}/\bar\beta$. Then for every configuration $\pi\in\Pi$ and
every unit $u$,
\begin{equation}\label{eq:Ceps}
(1-\varepsilon)\,C_u(\pi)\ \le\ \Chat_u(\pi)\ \le\ (1+\varepsilon)\,C_u(\pi),
\end{equation}
and the same two-sided multiplicative bound holds for every realized finish
time $\widehat f_u(\pi)$ versus $f_u(\pi)$ (work-dependent components; see
Remark~\ref{rem:latency} for latencies).
\end{lemma}

\begin{proof}
By Definition~\ref{def:struct}, $C_u(\pi)=\langle\kappa_{x_u},d_u\rangle$ and
$\Chat_u(\pi)=\langle\kappa_{x_u},\widehat d_u\rangle$ with the \emph{same}
nonnegative rate vector $\kappa_{x_u}=\kappa_{x_u}(\bar L_{x_u})$ (descriptors
are fixed; we estimate only demands). By \eqref{eq:Deps},
$(1-\varepsilon)d_u\le\widehat d_u\le(1+\varepsilon)d_u$ componentwise, so by
nonnegativity of $\kappa$,
\[
(1-\varepsilon)\langle\kappa,d_u\rangle\le\langle\kappa,\widehat d_u\rangle
\le(1+\varepsilon)\langle\kappa,d_u\rangle,
\]
which is \eqref{eq:Ceps}. Each transfer work term $d_{uv}/\bar\beta$ likewise
satisfies a two-sided $(1\pm\varepsilon)$ bound. By the left-shift recursion of
Lemma~\ref{lem:homog}, each finish time $f_u(\pi)$ is a finite max of sums of
such $(1\pm\varepsilon)$-bounded nonnegative terms; a max of sums of terms each
within a common multiplicative factor $(1\pm\varepsilon)$ is itself within the
same factor, giving the finish-time bound.
\end{proof}

The decisive gain over direct probing: \eqref{eq:Ceps} holds for \emph{all}
$\pi\in\Pi$ \emph{simultaneously} from the \emph{single} demand-accuracy event
$E_\varepsilon$, because the structural model carries the configuration
dependence deterministically. The sample budget is $O(Q\,\varepsilon_0^{-2}
\ln(Q/\delta))$, polynomial in the input size, with no dependence on $|\Pi|$.

\begin{remark}[Temporal drift]\label{rem:drift}
Stationarity over the probing horizon is needed for Assumption~\ref{ass:probe}.
If instead the demand/rate parameters drift Lipschitz-in-time with constant
$\eta$, i.e.\ $|d_\ell(\tau)-d_\ell(\tau')|\le\eta|\tau-\tau'|$, and the horizon
spans $H$ time units, then the additional systematic error is at most $\eta H$,
which is absorbed into the bias budget by replacing $\beta$ with
$\beta'=\beta+\eta H/a_{\min}$ in Proposition~\ref{prop:samples}; requiring
$\eta H\le\varepsilon'' a_{\min}$ keeps $\varepsilon''$ inside the budget. Thus
drift is handled by a horizon constraint folded into the same accounting.
\end{remark}

\section{Transfer theorem: accuracy $\Rightarrow$ mapping quality (part (c))}\label{sec:transfer}

We argue \emph{deterministically} on $E_\varepsilon$; the high-probability
statement follows by Corollary~\ref{cor:uniform}.

For a demand-scale-coherent $F$ define the \emph{estimated objective}
$\widehat F(\pi):=F(\widehat{\mathbf d},\pi)$ — i.e.\ $F$ evaluated with the
estimated finish times $\widehat f_u(\pi)$ — and the \emph{exact objective}
$F(\pi):=F(\mathbf d,\pi)$.

\begin{lemma}[Pointwise objective sandwich]\label{lem:sandwich}
Let $F$ be demand-scale-coherent and suppose $E_\varepsilon$ holds,
$\varepsilon\in(0,1)$. Then for every configuration $\pi\in\Pi$,
\begin{equation}\label{eq:sandwich}
(1-\varepsilon)\,F(\pi)\ \le\ \widehat F(\pi)\ \le\ (1+\varepsilon)\,F(\pi).
\end{equation}
\end{lemma}

\begin{proof}
By Lemma~\ref{lem:propagate}, $\widehat f_u(\pi)\le(1+\varepsilon)f_u(\pi)$ for
all $u$. By monotonicity of $F$ in the finish times and then degree-one
homogeneity (Definition~\ref{def:coherent}),
\[
\widehat F(\pi)=F\big((\widehat f_u(\pi))_u\big)\le F\big(((1+\varepsilon)f_u(\pi))_u\big)
=(1+\varepsilon)F(\pi).
\]
Symmetrically $\widehat f_u(\pi)\ge(1-\varepsilon)f_u(\pi)$ and $1-\varepsilon>0$
give $\widehat F(\pi)\ge(1-\varepsilon)F(\pi)$.
\end{proof}

\begin{theorem}[Approximation guarantee under estimated optimization]\label{thm:main}
Let $F$ be demand-scale-coherent and suppose $E_\varepsilon$ holds with
$\varepsilon\in(0,1)$. Let
$\hat\pi\in\arg\min_{\pi\in\Pi}\widehat F(\pi)$ and
$\pi^\star\in\arg\min_{\pi\in\Pi}F(\pi)$. Then
\begin{equation}\label{eq:approx}
F(\hat\pi)\ \le\ \frac{1+\varepsilon}{1-\varepsilon}\;F(\pi^\star).
\end{equation}
\end{theorem}

\begin{proof}
Both arg-min sets are nonempty since $\Pi$ is finite, nonempty. By the left
inequality of \eqref{eq:sandwich} at $\hat\pi$, optimality of $\hat\pi$ for
$\widehat F$ (so $\widehat F(\hat\pi)\le\widehat F(\pi^\star)$), and the right
inequality at $\pi^\star$,
\[
(1-\varepsilon)F(\hat\pi)\le\widehat F(\hat\pi)\le\widehat F(\pi^\star)
\le(1+\varepsilon)F(\pi^\star).
\]
Divide by $1-\varepsilon>0$.
\end{proof}

\begin{theorem}[Robustness to approximate optimization]\label{thm:approxalg}
Let $F$ be demand-scale-coherent, $E_\varepsilon$ hold, $\varepsilon\in(0,1)$,
and $\gamma\ge1$. If $\tilde\pi\in\Pi$ satisfies
$\widehat F(\tilde\pi)\le\gamma\min_{\pi}\widehat F(\pi)$, then
\begin{equation}\label{eq:approxalg}
F(\tilde\pi)\ \le\ \gamma\,\frac{1+\varepsilon}{1-\varepsilon}\;\min_{\pi\in\Pi}F(\pi).
\end{equation}
\end{theorem}

\begin{proof}
With $\pi^\star\in\arg\min_\pi F(\pi)$, chain the left inequality of
\eqref{eq:sandwich} at $\tilde\pi$, the $\gamma$-approximation hypothesis,
$\min_\pi\widehat F\le\widehat F(\pi^\star)$, and the right inequality at
$\pi^\star$:
\[
(1-\varepsilon)F(\tilde\pi)\le\widehat F(\tilde\pi)\le\gamma\min_\pi\widehat F
\le\gamma\widehat F(\pi^\star)\le\gamma(1+\varepsilon)F(\pi^\star).
\]
Divide by $1-\varepsilon$. Taking $\gamma=1$ recovers Theorem~\ref{thm:main}.
\end{proof}

This is the relevant form in practice: optimal DAG scheduling is NP-hard, but a
polynomial $\gamma$-approximation (e.g.\ Graham list scheduling with
$\gamma\le 2-\tfrac1n$ for makespan on identical machines, or its heterogeneous
variants) run \emph{on the estimated demands} yields a configuration within
$\gamma\frac{1+\varepsilon}{1-\varepsilon}$ of the true optimum.

\begin{remark}[Additive latency correction]\label{rem:latency}
If nonzero pure latencies $\lambda_{ij}>0$ appear, write
$f_u(\pi)=f_u^{\mathrm{work}}(\pi)+\Lambda_u(\pi)$ where $\Lambda_u(\pi)$ is the
total accumulated latency along the critical chain to $u$ (a
$\mathbf d$-independent, exactly known constant for fixed $\pi$). Then
$\widehat f_u(\pi)$ satisfies $|\widehat f_u-f_u|\le\varepsilon f_u^{\mathrm{work}}
\le\varepsilon f_u$, so \eqref{eq:sandwich} holds verbatim; if instead one wants
the latency excluded from the relative budget, replace $\varepsilon$ by
$\varepsilon\,\sup_\pi f_u^{\mathrm{work}}(\pi)/f_u(\pi)\le\varepsilon$ — the
bound only improves. Hence latencies never worsen the ratio.
\end{remark}

\begin{corollary}[End-to-end high-probability guarantee]\label{cor:e2e}
Fix $\varepsilon\in(0,1)$, $\beta\in[0,\varepsilon)$, $\delta\in(0,1)$, a
demand-scale-coherent $F$, and a $\gamma$-approximation oracle for
$\min_\pi\widehat F$. Run \textsc{ComputeDescriptors}, estimate
$\widehat{\mathbf d}$ with
$r_\ell^{\mathrm{eff}}\ge\frac{b_\ell^2}{2a_\ell^2(\varepsilon-\beta)^2}\ln\frac{2Q}{\delta}$
probes per coordinate (Assumption~\ref{ass:probe}), and output
$\tilde\pi$ from the oracle on $\widehat F$. Then with probability $\ge1-\delta$,
\[
F(\tilde\pi)\ \le\ \gamma\,\frac{1+\varepsilon}{1-\varepsilon}\,\min_{\pi\in\Pi}F(\pi).
\]
\end{corollary}

\begin{proof}
By Corollary~\ref{cor:uniform}, $E_\varepsilon$ holds w.p.\ $\ge1-\delta$; on it,
Lemma~\ref{lem:propagate} gives uniform completion/finish-time accuracy over all
$\pi$, and Theorem~\ref{thm:approxalg} gives the bound.
\end{proof}

\begin{remark}[Tightness]\label{rem:tight}
The factor $\frac{1+\varepsilon}{1-\varepsilon}$ is unimprovable for
demand-scale-coherent objectives. Take a single unit $u$, two machines, $F=F_{\max}$
(so $f_u=C_u$, no precedence), one demand coordinate $d$, rates
$\kappa_1=1-\varepsilon$, $\kappa_2=1+\varepsilon$, exact $d=1$, so true times
$C_{u1}=1-\varepsilon$, $C_{u2}=1+\varepsilon$. Suppose the estimate
$\widehat d$ realizes $\widehat C_{u1}=(1+\varepsilon)(1-\varepsilon)$ on machine
1 and $\widehat C_{u2}=(1-\varepsilon)(1+\varepsilon)$ on machine 2 (e.g.\ by the
model evaluating $\widehat d=(1+\varepsilon)$ for the comparison to machine 1 and
$(1-\varepsilon)$ for machine 2 in a worst-case multi-coordinate instance). The
two estimates coincide, so the optimizer may pick machine 1, giving
$F(\hat\pi)=1+\varepsilon$ while $F(\pi^\star)=1-\varepsilon$; both estimates obey
\eqref{eq:Deps}, so the ratio $\frac{1+\varepsilon}{1-\varepsilon}$ is attained.
\end{remark}

\begin{algorithm}[h]
\caption{EstimatedOptimalMapping}\label{alg:map}
\begin{algorithmic}[1]
\Require $\varepsilon\in(0,1)$, bias bound $\beta\in[0,\varepsilon)$, $\delta\in(0,1)$, demand-scale-coherent $F$, bounds $a_\ell,b_\ell$, autocorr.\ time $\tau_{\mathrm{ac}}$, $\gamma$-oracle, configuration family $\Pi$.
\Ensure $\tilde\pi$ with $F(\tilde\pi)\le\gamma\frac{1+\varepsilon}{1-\varepsilon}\min_\pi F(\pi)$ w.p.\ $\ge1-\delta$.
\State $(\bar L,\bar N)\gets$ \textsc{ComputeDescriptors}; set rates $\kappa_i(\bar L_i)$
\For{each demand coordinate $\ell=1,\dots,Q$}
  \State $r_\ell\gets\big\lceil \tau_{\mathrm{ac}}\frac{b_\ell^2}{2a_\ell^2(\varepsilon-\beta)^2}\ln\frac{2Q}{\delta}\big\rceil$
  \State draw probes $Y^{(1)}_\ell,\dots,Y^{(r_\ell)}_\ell$ per Assumption~\ref{ass:probe}; $\widehat d_\ell\gets\frac1{r_\ell}\sum_p Y^{(p)}_\ell$
\EndFor
\State define $\widehat F(\pi):=F(\widehat{\mathbf d},\pi)$ via the structural model and left-shift recursion
\State $\tilde\pi\gets\gamma$-oracle$\big(\min_{\pi\in\Pi}\widehat F(\pi)\big)$
\State \Return $\tilde\pi$
\end{algorithmic}
\end{algorithm}

\noindent\textbf{Correctness.} Lines 2--5 realize the demand estimator at the
sample size of Corollary~\ref{cor:uniform}, so $E_\varepsilon$ holds w.p.\
$\ge1-\delta$; lines 6--7 evaluate $\widehat F$ deterministically through the
structural model (valid for \emph{all} $\pi$ by Lemma~\ref{lem:propagate}) and
return a $\gamma$-approximate minimizer. Corollary~\ref{cor:e2e} certifies the
output guarantee. \textbf{Termination:} finite loops, finite $\Pi$.
\textbf{Complexity:} probing costs
$\sum_\ell r_\ell=O\!\big(\tau_{\mathrm{ac}}\,Q\,\frac{b_{\max}^2}{a_{\min}^2(\varepsilon-\beta)^2}\ln\frac{Q}{\delta}\big)$
probes — polynomial in the instance and \emph{independent of $|\Pi|$}; line~7
costs one call to the scheduling oracle. The reduction is black-box in the
optimizer, so the heterogeneous scheduling literature's approximation algorithms
plug in directly.

\section{Summary of what is established}\label{sec:summary}

(a) Descriptors $L_i(\tau),N_{ij}(\tau)$ are defined and computed by
\textsc{ComputeDescriptors} (window averaging), feeding the structural rate
vectors. (b) The estimator $\widehat{\mathbf d}$ of demands, propagated through
the known structural completion-time model, yields $\Chat$ and finish times that
are uniformly $(1\pm\varepsilon)$-accurate \emph{over all configurations}, with
explicit sample complexity that honestly carries bias ($\beta$), sample
dependence ($\tau_{\mathrm{ac}}$), and temporal drift ($\eta H$). (c) The
transfer theorem (Theorems~\ref{thm:main}, \ref{thm:approxalg},
Corollary~\ref{cor:e2e}) shows that minimizing the estimated objective — exactly
or $\gamma$-approximately — yields a configuration within
$\gamma\frac{1+\varepsilon}{1-\varepsilon}$ of the true optimum for every
demand-scale-coherent objective (makespan under precedence, weighted completion
time, $\ell_p$ loads, $\dots$), and the constant $\frac{1+\varepsilon}{1-\varepsilon}$
is tight.
\end{document}
